% !TeX root = ../../main.tex
% =====================================================================
%  El coche como sistema
%  Responsable:
%  Última revisión:
%  NOTA: capítulo deliberadamente breve. Describe el método al que el equipo
%  quiere llegar, no el que ha seguido esta temporada. Ver el recuadro CASO.
% =====================================================================
\chapter{El coche como sistema}
\label{ch:coche_como_sistema}

\begin{objetivos}
  \item Distinguir un objetivo de vehículo de un deseo.
  \item Explicar cómo un objetivo global se convierte en un requisito concreto de
        una pieza.
  \item Entender qué es un presupuesto de masa y para qué sirve un modelo de
        centro de gravedad.
  \item Explicar por qué la simulación de vuelta es el único árbitro razonable
        entre subsistemas que compiten por lo mismo.
\end{objetivos}

\fichasesion{60 minutos}{Capítulo \ref{ch:formula_student}}%
{Hoja de cálculo de masas del coche actual, si existe}

Este capítulo es corto a propósito, y conviene decir por qué desde el principio:
describe la forma de trabajar a la que el equipo quiere llegar, no la que ha
seguido hasta ahora. Se incluye en el tronco común porque es el marco que da
sentido a los capítulos de dinámica vehicular, aerodinámica y transmisión, y
porque los jueces del evento de diseño preguntan por él antes que por cualquier
detalle de una pieza.

La idea de fondo es simple. Un coche de competición es un sistema en el que casi
ningún subsistema puede mejorar sin empeorar otro: el aerodinámico añade carga y
también masa y resistencia; el radiador necesita aire que la carrocería quería
para otra cosa; una relación de transmisión corta gana en aceleración y pierde en
velocidad punta. Si cada subsistema optimiza lo suyo por separado, el resultado no
es un coche optimizado, sino un coche incoherente. Hace falta un criterio común
--- unos objetivos de vehículo --- y una herramienta que traduzca cualquier
decisión local a ese criterio.

% ---------------------------------------------------------------------
\section{Objetivos de vehículo}
\label{sec:coche_como_sistema:01}
\index{objetivos de vehículo}

Un objetivo de vehículo es un número, con fecha y con responsable. «El coche tiene
que ser ligero» no es un objetivo; «masa en orden de marcha sin piloto por debajo
de \SI{230}{\kilogram}, verificada en báscula antes del 15 de mayo, responsable
el jefe de chasis» sí lo es.

Los objetivos típicos de un coche de Formula Student son media docena:

\begin{center}
\begin{tabularx}{\textwidth}{@{}lXl@{}}
\toprule
\textbf{Objetivo} & \textbf{Por qué se fija} & \textbf{Unidad} \\
\midrule
Masa total & Afecta a todo: aceleración, frenada, curva, consumo & \si{\kilogram} \\
Reparto de masas & Condiciona el equilibrio y la tracción & \si{\percent} \\
Altura del centro de gravedad & Gobierna la transferencia de carga & \si{\milli\meter} \\
Rigidez torsional del chasis & Determina si la suspensión hace lo que se diseñó & \si{\newton\meter\per\degree} \\
Carga aerodinámica y resistencia & Prestaciones en curva frente a recta & \si{\meter\squared} \\
Tiempo objetivo por vuelta & Integra todo lo anterior en un solo número & \si{\second} \\
\bottomrule
\end{tabularx}
\end{center}

Lo importante no es acertar el valor a la primera --- casi nunca se acierta ---
sino que exista, que esté escrito y que cuando se incumpla \emph{se sepa}. Si un
objetivo se revisa a la baja, la revisión se discute y se anota.

\begin{clave}
El valor de un objetivo de vehículo no está en el número, sino en que obliga a
tener la discusión \emph{antes} de fabricar. Cuando dos subsistemas piden lo
mismo, el objetivo es lo que permite decidir sin que gane el que más levanta la voz.
\end{clave}

% ---------------------------------------------------------------------
\section{Cascada de requisitos}
\label{sec:coche_como_sistema:02}
\index{cascada de requisitos}

Un objetivo de vehículo no se puede diseñar directamente: hay que descomponerlo
hasta llegar a algo que una persona pueda dimensionar. Ese descenso es la cascada
de requisitos, y tiene tres niveles:

\begin{enumerate}
  \item \textbf{Objetivo de vehículo.} Masa total por debajo de
        \SI{230}{\kilogram}.
  \item \textbf{Requisito de subsistema.} Suspensión completa (cuatro esquinas)
        por debajo de \SI{28}{\kilogram}.
  \item \textbf{Requisito de componente.} Trapecio delantero superior por debajo
        de \SI{380}{\gram} cumpliendo los casos de carga del capítulo
        \ref{ch:casos_carga}.
\end{enumerate}

La cascada funciona en los dos sentidos. Hacia abajo reparte el presupuesto; hacia
arriba avisa. Si el trapecio se va a \SI{450}{\gram} y no hay forma de bajarlo,
alguien tiene que decidir conscientemente si se acepta, si se compensa en otro
sitio o si se cambia el objetivo. Lo que no puede pasar es que nadie se entere
hasta la báscula.

% ---------------------------------------------------------------------
\section{Presupuesto de masa y centro de gravedad}
\label{sec:coche_como_sistema:03}
\index{presupuesto de masa}\index{centro de gravedad}

El presupuesto de masa es una hoja de cálculo con una fila por pieza y tres
columnas de masa: \textbf{estimada} (al fijar el objetivo), \textbf{de CAD}
(cuando la pieza está modelada, con el material asignado de verdad) y
\textbf{pesada} (cuando existe). La diferencia entre la primera y la última
columna indica cómo de fiables son nuestras estimaciones, y sirve para ajustarlas
la temporada siguiente.

Si además se anota la posición del centro de gravedad de cada pieza, la misma
hoja da el centro de gravedad del coche completo:

\begin{equation}
  x_{\CdG} = \frac{\sum_i m_i\, x_i}{\sum_i m_i}
  \label{eq:cdg}
\end{equation}

y lo mismo para $y$ y $z$. No hace falta nada más sofisticado. Con eso se puede
responder a preguntas del tipo «¿cuánto sube el centro de gravedad si monto el
radiador aquí en vez de allí?» en el momento en que se plantean, que es cuando
sirve de algo.

La altura del centro de gravedad merece atención especial porque gobierna la
transferencia de carga (capítulo \ref{ch:regimen_estacionario}) y, con ella, el
agarre disponible. Suele recibir mucha menos atención que la masa total, y en un
coche de Formula Student unos pocos centímetros de altura pesan tanto como varios
kilogramos.

\begin{ojo}
\begin{itemize}
  \item Actualizar el presupuesto de masa solo un par de veces al año: deja de
        servir para decidir. Se actualiza cada vez que se libera una pieza.
  \item Estimar masas «por experiencia» sin apuntar de dónde sale la estimación:
        cuando falla, no se puede aprender nada.
  \item Olvidar en la hoja lo que no es una pieza: cableado, líquidos, tornillería,
        pintura, cinta. Suele ser bastante más de lo que nadie espera.
  \item Medir el centro de gravedad solo en el plano y no en altura, que es lo
        que realmente condiciona el comportamiento.
\end{itemize}
\end{ojo}

% ---------------------------------------------------------------------
\section{La simulación de vuelta como árbitro}
\label{sec:coche_como_sistema:04}
\index{simulación de vuelta}

Supongamos la discusión clásica: el paquete aerodinámico añade
\SI{0.9}{\meter\squared} de $\ClA$ y \SI{12}{\kilogram} de masa. ¿Compensa? No
hay forma de contestar razonando: hay que calcularlo, porque una parte del efecto
es favorable y otra no, y el resultado depende del trazado, de las velocidades y
del agarre disponible.

Eso es exactamente lo que hace una simulación de vuelta: convierte un cambio de
diseño en segundos por vuelta y, a través de la tabla de puntos del capítulo
\ref{ch:formula_student}, en puntos. Su valor real no está en predecir el tiempo
absoluto --- eso casi nunca lo acierta --- sino en dar \textbf{sensibilidades}:

\begin{center}
\begin{tabular}{@{}ll@{}}
\toprule
\textbf{Sensibilidad} & \textbf{Unidades} \\
\midrule
Segundos por vuelta por kilogramo & \si{\second\per\kilogram} \\
Segundos por vuelta por unidad de $\ClA$ & \si{\second\per\meter\squared} \\
Segundos por vuelta por milímetro de altura de $\CdG$ & \si{\second\per\milli\meter} \\
Segundos por vuelta por punto de coeficiente de agarre & \si{\second} \\
\bottomrule
\end{tabular}
\end{center}

Con esa tabla, la discusión anterior se resuelve con un número. Y, de cara al
evento de diseño, la respuesta a «¿por qué lleváis alerón?» deja de ser una
opinión.

Estas sensibilidades son la entrada de los capítulos
\ref{ch:merece_la_pena_aero} (aerodinámica) y \ref{ch:relacion_transmision}
(relación de transmisión). Sin ellas, esos dos capítulos se quedan en criterio
cualitativo.

\begin{clave}
La simulación de vuelta no sirve para saber cuánto se tarda: sirve para saber
\textbf{cuánto cambia el tiempo} cuando cambias algo. Un modelo mediocre bien
entendido da sensibilidades útiles; un modelo excelente sin validar da un número
absoluto en el que nadie debería confiar.
\end{clave}

% ---------------------------------------------------------------------
\section{Interfaces entre subsistemas}
\label{sec:coche_como_sistema:05}
\index{interfaces}

Los fallos de integración no ocurren dentro de los subsistemas, sino en las
fronteras, y aparecen tarde: en el montaje. Se previenen con una tabla de
interfaces que se acuerda al principio de la temporada y en la que cada línea
tiene un dueño único.

\begin{center}
\begin{tabularx}{\textwidth}{@{}llX@{}}
\toprule
\textbf{Interfaz} & \textbf{Dueño} & \textbf{Qué hay que acordar} \\
\midrule
Chasis--suspensión & Chasis & Posición y tolerancia de los anclajes \\
Suspensión--freno--rueda & Suspensión & Espacio libre en la llanta, cotas de buje \\
Refrigeración--aerodinámica & Aerodinámica & Caudal y superficie de entrada de aire \\
Chasis--transmisión & Transmisión & Anclajes, alineación y resguardos \\
Chasis--ergonomía & Ergonomía & Volumen de cockpit, plantillas de reglamento \\
Electrónica--todos & Electrónica & Recorrido de mazos, sensores, masas \\
\bottomrule
\end{tabularx}
\end{center}

La regla práctica es que \textbf{ninguna interfaz puede tener dos dueños}. Si dos
subsistemas comparten una cota, uno la fija y el otro la consume; y la cota vive
en el esqueleto del CAD (capítulo \ref{ch:herramientas_cad_datos}), no en la
cabeza de nadie.

% ---------------------------------------------------------------------
%  Cierre del capítulo
% ---------------------------------------------------------------------

\begin{caso}[dónde estamos de verdad]
Esta temporada el equipo \textbf{no ha trabajado con objetivos de vehículo
formales}: no se fijó un objetivo de masa, ni de altura de centro de gravedad, ni
un tiempo por vuelta de referencia, y las decisiones de cada subsistema se han
tomado con criterio local. Tampoco se ha usado una simulación de vuelta para
arbitrar entre subsistemas.

Conviene escribirlo tal cual, por dos razones. La primera es que este libro solo
sirve si es honesto. La segunda es que \emph{es exactamente la pregunta que hacen
los jueces de diseño}: cómo se decidió el reparto de prioridades del coche. Tener
clara la respuesta ---incluso cuando la respuesta es «este año no lo hicimos
así, y estas son las consecuencias que hemos visto»--- puntúa más que improvisar
una justificación.

El camino para la próxima temporada es el que describe este capítulo, en este orden:

\begin{enumerate}
  \item Montar el presupuesto de masa con lo que ya existe, aunque sea a
        posteriori: pesar el coche actual pieza a pieza. Está comprometido y
        pendiente de hacer; es la base de todo lo demás.
  \item Medir el centro de gravedad del coche actual, en las tres direcciones.
  \item Levantar una simulación de vuelta sencilla, aunque sea cuasi-estacionaria
        (capítulo \ref{ch:simulacion_vuelta}), y sacar las cuatro sensibilidades
        de la tabla anterior.
  \item Con esos datos ya sí: fijar los objetivos de la temporada siguiente.
\end{enumerate}

\redactar{Actualizar este recuadro cuando se completen los pasos anteriores,
dejando constancia de las cifras reales obtenidas.}
\end{caso}

\begin{ojo}
\begin{itemize}
  \item Fijar objetivos de vehículo copiando los de otro equipo. Un objetivo que
        el equipo no sabe de dónde sale no se defiende ante un juez ni se cumple.
  \item Poner objetivos tan holgados que se cumplen solos: no informan de nada.
  \item Usar la simulación de vuelta para decidir y no validarla nunca contra
        datos de pista.
  \item Descubrir las interfaces en el montaje.
\end{itemize}
\end{ojo}

\begin{entregable}
\begin{enumerate}
  \item La hoja de presupuesto de masa del coche actual, con las columnas de masa
        estimada, de CAD y pesada.
  \item La tabla de interfaces con un dueño por línea, acordada y firmada por los
        responsables de subsistema.
\end{enumerate}
\end{entregable}

\begin{juez}
\begin{enumerate}
  \item ¿Qué objetivos de vehículo os fijasteis y de dónde salieron esos valores?
  \item ¿Cuánto pesa el coche y cuánto pensabais que iba a pesar al empezar a
        diseñarlo?
  \item ¿Dónde está el centro de gravedad y cómo lo habéis medido?
  \item ¿Cuántos segundos por vuelta os cuesta cada kilogramo?
  \item Cuando aerodinámica y refrigeración han pedido el mismo espacio, ¿quién
        ha decidido y con qué criterio?
\end{enumerate}
\end{juez}
